\section*{NeurIPS Paper Checklist}

\begin{enumerate}

\item {\bf Claims}
    \item[] Question: Do the main claims made in the abstract and introduction accurately reflect the paper's contributions and scope?
    \item[] Answer: \answerYes{}
    \item[] Justification: The abstract and introduction explicitly enumerate the five contributions (structural role separation, verifier as binding constraint, cross-provider role mixing, prompt-only versus enforced, open release of all artifacts). 
    The numerical claims in the abstract are tied to corresponding sample sizes, confidence intervals where applicable, and detailed tables in \S\ref{sec:experiment} and the appendix.
    \item[] Guidelines:
    \begin{itemize}
        \item The answer \answerNA{} means that the abstract and introduction do not include the claims made in the paper.
        \item The abstract and/or introduction should clearly state the claims made, including the contributions made in the paper and important assumptions and limitations. A \answerNo{} or \answerNA{} answer to this question will not be perceived well by the reviewers. 
        \item The claims made should match theoretical and experimental results, and reflect how much the results can be expected to generalize to other settings. 
        \item It is fine to include aspirational goals as motivation as long as it is clear that these goals are not attained by the paper. 
    \end{itemize}

\item {\bf Limitations}
    \item[] Question: Does the paper discuss the limitations of the work performed by the authors?
    \item[] Answer: \answerYes{}
    \item[] Justification: \S\ref{sec:errors} reports the open-source tool-call reliability limit, \S\ref{sec:promptonly} reports power-limited prompt-only contrasts and treats the shared-history condition as inconclusive, and the Conclusion lists deferred axes (multi-round dialogue, dynamic role assignment, within-provider scaling). Appendix~\ref{app:enforcement-sensitivity} reports the worst-case sensitivity treatment.
    \item[] Guidelines:
    \begin{itemize}
        \item The answer \answerNA{} means that the paper has no limitation while the answer \answerNo{} means that the paper has limitations, but those are not discussed in the paper. 
        \item The authors are encouraged to create a separate ``Limitations'' section in their paper.
        \item The paper should point out any strong assumptions and how robust the results are to violations of these assumptions (e.g., independence assumptions, noiseless settings, model well-specification, asymptotic approximations only holding locally). The authors should reflect on how these assumptions might be violated in practice and what the implications would be.
        \item The authors should reflect on the scope of the claims made, e.g., if the approach was only tested on a few datasets or with a few runs. In general, empirical results often depend on implicit assumptions, which should be articulated.
        \item The authors should reflect on the factors that influence the performance of the approach. For example, a facial recognition algorithm may perform poorly when image resolution is low or images are taken in low lighting. Or a speech-to-text system might not be used reliably to provide closed captions for online lectures because it fails to handle technical jargon.
        \item The authors should discuss the computational efficiency of the proposed algorithms and how they scale with dataset size.
        \item If applicable, the authors should discuss possible limitations of their approach to address problems of privacy and fairness.
        \item While the authors might fear that complete honesty about limitations might be used by reviewers as grounds for rejection, a worse outcome might be that reviewers discover limitations that aren't acknowledged in the paper. The authors should use their best judgment and recognize that individual actions in favor of transparency play an important role in developing norms that preserve the integrity of the community. Reviewers will be specifically instructed to not penalize honesty concerning limitations.
    \end{itemize}

\item {\bf Theory assumptions and proofs}
    \item[] Question: For each theoretical result, does the paper provide the full set of assumptions and a complete (and correct) proof?
    \item[] Answer: \answerNA{}
    \item[] Justification: The paper introduces a benchmark and a controlled experimental design. There are no theoretical results.
    \item[] Guidelines:
    \begin{itemize}
        \item The answer \answerNA{} means that the paper does not include theoretical results. 
        \item All the theorems, formulas, and proofs in the paper should be numbered and cross-referenced.
        \item All assumptions should be clearly stated or referenced in the statement of any theorems.
        \item The proofs can either appear in the main paper or the supplemental material, but if they appear in the supplemental material, the authors are encouraged to provide a short proof sketch to provide intuition. 
        \item Inversely, any informal proof provided in the core of the paper should be complemented by formal proofs provided in appendix or supplemental material.
        \item Theorems and Lemmas that the proof relies upon should be properly referenced. 
    \end{itemize}

\item {\bf Experimental result reproducibility}
    \item[] Question: Does the paper fully disclose all the information needed to reproduce the main experimental results of the paper to the extent that it affects the main claims and/or conclusions of the paper?
    \item[] Answer: \answerYes{}
    \item[] Justification: Appendix~\ref{app:reproducibility} documents the model list, prompts, dropped task identifiers, replication run counts (1{,}165 reference-ablation, 2{,}025 role-mixing), seed regime, and total spend. All graders are deterministic shell scripts and all task generators are deterministic given a seed. The benchmark, generators, and harness are released under a permissive license.
    \item[] Guidelines:
    \begin{itemize}
        \item The answer \answerNA{} means that the paper does not include experiments.
        \item If the paper includes experiments, a \answerNo{} answer to this question will not be perceived well by the reviewers: Making the paper reproducible is important, regardless of whether the code and data are provided or not.
        \item If the contribution is a dataset and\slash or model, the authors should describe the steps taken to make their results reproducible or verifiable. 
        \item Depending on the contribution, reproducibility can be accomplished in various ways. For example, if the contribution is a novel architecture, describing the architecture fully might suffice, or if the contribution is a specific model and empirical evaluation, it may be necessary to either make it possible for others to replicate the model with the same dataset, or provide access to the model. In general. releasing code and data is often one good way to accomplish this, but reproducibility can also be provided via detailed instructions for how to replicate the results, access to a hosted model (e.g., in the case of a large language model), releasing of a model checkpoint, or other means that are appropriate to the research performed.
        \item While NeurIPS does not require releasing code, the conference does require all submissions to provide some reasonable avenue for reproducibility, which may depend on the nature of the contribution. For example
        \begin{enumerate}
            \item If the contribution is primarily a new algorithm, the paper should make it clear how to reproduce that algorithm.
            \item If the contribution is primarily a new model architecture, the paper should describe the architecture clearly and fully.
            \item If the contribution is a new model (e.g., a large language model), then there should either be a way to access this model for reproducing the results or a way to reproduce the model (e.g., with an open-source dataset or instructions for how to construct the dataset).
            \item We recognize that reproducibility may be tricky in some cases, in which case authors are welcome to describe the particular way they provide for reproducibility. In the case of closed-source models, it may be that access to the model is limited in some way (e.g., to registered users), but it should be possible for other researchers to have some path to reproducing or verifying the results.
        \end{enumerate}
    \end{itemize}

\item {\bf Open access to data and code}
    \item[] Question: Does the paper provide open access to the data and code, with sufficient instructions to faithfully reproduce the main experimental results?
    \item[] Answer: \answerYes{}
    \item[] Justification: The benchmark, all 851 task templates with deterministic shell-script graders, the parameterized generators, the agent harness, the role-mixing protocol, the 5-condition ablation runner, the audited \textsc{TeamBench-Verified} subset, and the human-baseline collection platform are released under a permissive open-source license at the project URL given in Appendix~\ref{app:reproducibility}. The leaderboard task selection, per-condition reference-ablation logs, deduplicated role-mixing per-task records, and cost ledger are included.
    \item[] Guidelines:
    \begin{itemize}
        \item The answer \answerNA{} means that paper does not include experiments requiring code.
        \item Please see the NeurIPS code and data submission guidelines (\url{https://neurips.cc/public/guides/CodeSubmissionPolicy}) for more details.
        \item While we encourage the release of code and data, we understand that this might not be possible, so \answerNo{} is an acceptable answer. Papers cannot be rejected simply for not including code, unless this is central to the contribution (e.g., for a new open-source benchmark).
        \item The instructions should contain the exact command and environment needed to run to reproduce the results. See the NeurIPS code and data submission guidelines (\url{https://neurips.cc/public/guides/CodeSubmissionPolicy}) for more details.
        \item The authors should provide instructions on data access and preparation, including how to access the raw data, preprocessed data, intermediate data, and generated data, etc.
        \item The authors should provide scripts to reproduce all experimental results for the new proposed method and baselines. If only a subset of experiments are reproducible, they should state which ones are omitted from the script and why.
        \item At submission time, to preserve anonymity, the authors should release anonymized versions (if applicable).
        \item Providing as much information as possible in supplemental material (appended to the paper) is recommended, but including URLs to data and code is permitted.
    \end{itemize}

\item {\bf Experimental setting/details}
    \item[] Question: Does the paper specify all the training and test details?
    \item[] Answer: \answerYes{}
    \item[] Justification: \S\ref{sec:exp-studies} documents step budgets ($15$/$25$/$10$ turns for Planner/Executor/Verifier), temperature ($0$), tool set, deterministic-grader rule (binary pass requires every check to pass), and statistical procedure (paired bootstrap, $10{,}000$ iterations, Wilson $95\%$ CIs, McNemar with Holm-Bonferroni for the enforcement study). Per-role tool access, file scoping, system prompts, and Docker bind-mount permissions are in Appendices~\ref{app:role-permissions}, \ref{app:tools}, \ref{app:prompts}.
    \item[] Guidelines:
    \begin{itemize}
        \item The answer \answerNA{} means that the paper does not include experiments.
        \item The experimental setting should be presented in the core of the paper to a level of detail that is necessary to appreciate the results and make sense of them.
        \item The full details can be provided either with the code, in appendix, or as supplemental material.
    \end{itemize}

\item {\bf Experiment statistical significance}
    \item[] Question: Does the paper report error bars suitably and correctly defined?
    \item[] Answer: \answerYes{}
    \item[] Justification: All cross-model comparisons use paired bootstrap ($10{,}000$ iterations) with Wilson $95\%$ CIs. Per-task and per-quintile uplift use paired-bootstrap CIs (Fig.~\ref{fig:error-analysis}(b)). The pre-registered enforcement study uses McNemar with Holm-Bonferroni correction over three planned tests, with both raw and adjusted $p$-values reported (Table~\ref{tab:promptonly}). Wilson CIs are reported in figure whiskers and the prompt-only results table.
    \item[] Guidelines:
    \begin{itemize}
        \item The answer \answerNA{} means that the paper does not include experiments.
        \item The authors should answer \answerYes{} if the results are accompanied by error bars, confidence intervals, or statistical significance tests, at least for the experiments that support the main claims of the paper.
        \item The factors of variability that the error bars are capturing should be clearly stated (for example, train/test split, initialization, random drawing of some parameter, or overall run with given experimental conditions).
        \item The method for calculating the error bars should be explained (closed form formula, call to a library function, bootstrap, etc.)
        \item The assumptions made should be given (e.g., Normally distributed errors).
        \item It should be clear whether the error bar is the standard deviation or the standard error of the mean.
        \item It is OK to report 1-sigma error bars, but one should state it. The authors should preferably report a 2-sigma error bar than state that they have a 96\% CI, if the hypothesis of Normality of errors is not verified.
        \item For asymmetric distributions, the authors should be careful not to show in tables or figures symmetric error bars that would yield results that are out of range (e.g., negative error rates).
        \item If error bars are reported in tables or plots, the authors should explain in the text how they were calculated and reference the corresponding figures or tables in the text.
    \end{itemize}

\item {\bf Experiments compute resources}
    \item[] Question: For each experiment, does the paper provide sufficient information on the computer resources?
    \item[] Answer: \answerYes{}
    \item[] Justification: Appendix~\ref{app:reproducibility} reports total measured spend (\$326.04 for the cross-provider grid, $1{,}165$ runs for the reference ablation under gemini-3-flash-preview). Inference is run via vendor APIs for commercial models and via vLLM on local GPU hosts for open-weight models. No fine-tuning or model training was performed.
    \item[] Guidelines:
    \begin{itemize}
        \item The answer \answerNA{} means that the paper does not include experiments.
        \item The paper should indicate the type of compute workers CPU or GPU, internal cluster, or cloud provider, including relevant memory and storage.
        \item The paper should provide the amount of compute required for each of the individual experimental runs as well as estimate the total compute. 
        \item The paper should disclose whether the full research project required more compute than the experiments reported in the paper (e.g., preliminary or failed experiments that didn't make it into the paper). 
    \end{itemize}

\item {\bf Code of ethics}
    \item[] Question: Does the research conform with the NeurIPS Code of Ethics?
    \item[] Answer: \answerYes{}
    \item[] Justification: The benchmark is open-source and intended for academic research. The human-baseline platform follows the IRB protocol referenced under item~14 below. Adversarial-trap and security-vulnerability tasks contain plausible-but-incorrect security patterns by design and are flagged with per-task usage notes. Recommended use is in network-isolated containers, per the responsible-use note in Appendix~\ref{app:reproducibility}.
    \item[] Guidelines:
    \begin{itemize}
        \item The answer \answerNA{} means that the authors have not reviewed the NeurIPS Code of Ethics.
        \item If the authors answer \answerNo, they should explain the special circumstances that require a deviation from the Code of Ethics.
        \item The authors should make sure to preserve anonymity (e.g., if there is a special consideration due to laws or regulations in their jurisdiction).
    \end{itemize}

\item {\bf Broader impacts}
    \item[] Question: Does the paper discuss both potential positive societal impacts and negative societal impacts?
    \item[] Answer: \answerYes{}
    \item[] Justification: The benchmark surfaces a $49\%$ Verifier false-accept ceiling that is a direct safety signal: any production system that delegates correctness adjudication to a current LLM Verifier inherits this rate. The cross-provider results enable lower-cost and lower-vendor-lock procurement. Negative impact: deployable insights into multi-agent failure modes could in principle be used to red-team production agent systems. We mitigate by releasing graders and traces under a permissive license, putting attackers and defenders on the same footing.
    \item[] Guidelines:
    \begin{itemize}
        \item The answer \answerNA{} means that there is no societal impact of the work performed.
        \item If the authors answer \answerNA{} or \answerNo, they should explain why their work has no societal impact or why the paper does not address societal impact.
        \item Examples of negative societal impacts include potential malicious or unintended uses (e.g., disinformation, generating fake profiles, surveillance), fairness considerations (e.g., deployment of technologies that could make decisions that unfairly impact specific groups), privacy considerations, and security considerations.
        \item The conference expects that many papers will be foundational research and not tied to particular applications, let alone deployments. However, if there is a direct path to any negative applications, the authors should point it out. For example, it is legitimate to point out that an improvement in the quality of generative models could be used to generate Deepfakes for disinformation. On the other hand, it is not needed to point out that a generic algorithm for optimizing neural networks could enable people to train models that generate Deepfakes faster.
        \item The authors should consider possible harms that could arise when the technology is being used as intended and functioning correctly, harms that could arise when the technology is being used as intended but gives incorrect results, and harms following from (intentional or unintentional) misuse of the technology.
        \item If there are negative societal impacts, the authors could also discuss possible mitigation strategies (e.g., gated release of models, providing defenses in addition to attacks, mechanisms for monitoring misuse, mechanisms to monitor how a system learns from feedback over time, improving the efficiency and accessibility of ML).
    \end{itemize}

\item {\bf Safeguards}
    \item[] Question: Does the paper describe safeguards for responsible release?
    \item[] Answer: \answerYes{}
    \item[] Justification: Adversarial-trap (\textsc{TRAP}) and cryptographic (\textsc{CRYPTO}) tasks contain intentional weaknesses (nonce reuse, low PBKDF2 iterations, truncated authentication tags). Each task carries a metadata flag for adversarial content. Appendix~\ref{app:reproducibility} explicitly recommends running hosted leaderboard submissions in network-isolated containers. The benchmark releases no model weights and no scraped personal data.
    \item[] Guidelines:
    \begin{itemize}
        \item The answer \answerNA{} means that the paper poses no such risks.
        \item Released models that have a high risk for misuse or dual-use should be released with necessary safeguards to allow for controlled use of the model, for example by requiring that users adhere to usage guidelines or restrictions to access the model or implementing safety filters. 
        \item Datasets that have been scraped from the Internet could pose safety risks. The authors should describe how they avoided releasing unsafe images.
        \item We recognize that providing effective safeguards is challenging, and many papers do not require this, but we encourage authors to take this into account and make a best faith effort.
    \end{itemize}

\item {\bf Licenses for existing assets}
    \item[] Question: Are creators of existing assets properly credited and are licenses properly respected?
    \item[] Answer: \answerYes{}
    \item[] Justification: GitHub-sourced tasks credit upstream repositories (Flask, Click, httpx, Requests, Pydantic, Django, pytest, FastAPI, SQLAlchemy, Celery, Werkzeug, NumPy, SciPy, Keras, spaCy) and use only public-issue text under each project's license. Real data-science tasks use canonical UCI datasets~\citep{dua2019uci}. Real incident response tasks adapt public post-mortems. All commercial model APIs (Anthropic, Google, OpenAI) are accessed under their published terms of service. Open-weight model usage (Gemma, Qwen, gpt-oss) follows each model's stated license.
    \item[] Guidelines:
    \begin{itemize}
        \item The answer \answerNA{} means that the paper does not use existing assets.
        \item The authors should cite the original paper that produced the code package or dataset.
        \item The authors should state which version of the asset is used and, if possible, include a URL.
        \item The name of the license (e.g., CC-BY 4.0) should be included for each asset.
        \item For scraped data from a particular source (e.g., website), the copyright and terms of service of that source should be provided.
        \item If assets are released, the license, copyright information, and terms of use in the package should be provided. For popular datasets, \url{paperswithcode.com/datasets} has curated licenses for some datasets. Their licensing guide can help determine the license of a dataset.
        \item For existing datasets that are re-packaged, both the original license and the license of the derived asset (if it has changed) should be provided.
        \item If this information is not available online, the authors are encouraged to reach out to the asset's creators.
    \end{itemize}

\item {\bf New assets}
    \item[] Question: Are new assets introduced in the paper well documented?
    \item[] Answer: \answerYes{}
    \item[] Justification: The benchmark, harness, generators, ablation runner, role-mixing protocol, audit pillars, human-baseline collection platform, and pre-registered analysis plans are all documented in the appendix and accompanied by a dataset card following \citet{gebru2021datasheets}. Per-task usage notes flag adversarial content. The repository at the project URL contains a versioned release with an MIT license.
    \item[] Guidelines:
    \begin{itemize}
        \item The answer \answerNA{} means that the paper does not release new assets.
        \item Researchers should communicate the details of the dataset\slash code\slash model as part of their submissions via structured templates. This includes details about training, license, limitations, etc. 
        \item The paper should discuss whether and how consent was obtained from people whose asset is used.
        \item At submission time, remember to anonymize your assets (if applicable). You can either create an anonymized URL or include an anonymized zip file.
    \end{itemize}

\item {\bf Crowdsourcing and research with human subjects}
    \item[] Question: For research with human subjects, does the paper include the full text of instructions, screenshots, and details about compensation?
    \item[] Answer: \answerYes{}
    \item[] Justification: The Human Baseline collection platform is described in Appendix~\ref{app:human-platform} with the full four-screen flow (profile entry, task selection, mode/role selection, sandboxed workspace). The three survey instruments (Solo Reflection, Hybrid AI Teammate, CATME-lite + Coordination) are documented in the appendix with item lists. Pilot participants are unpaid academic volunteers and the platform supports per-task and per-role logging. The pre-registered analysis plan (Appendix~\ref{app:human-prereg}) commits eligibility filters, minimum-cell rules, and pairing strategy ahead of camera-ready analysis.
    \item[] Guidelines:
    \begin{itemize}
        \item The answer \answerNA{} means that the paper does not involve crowdsourcing nor research with human subjects.
        \item Including this information in the supplemental material is fine, but if the main contribution of the paper involves human subjects, then as much detail as possible should be included in the main paper. 
        \item According to the NeurIPS Code of Ethics, workers involved in data collection, curation, or other labor should be paid at least the minimum wage in the country of the data collector. 
    \end{itemize}

\item {\bf Institutional review board (IRB) approvals or equivalent for research with human subjects}
    \item[] Question: Does the paper describe potential risks, disclosure, and IRB approval?
    \item[] Answer: \answerYes{}
    \item[] Justification: The human-baseline pilot is covered by an institutional IRB exempt determination under U.S.\ federal regulation 45 CFR 46.104(d)(2), Exempt Category 2 (Educational Testing, Surveys, Interviews or Observation). The covered period spans the pilot reported here. Risk to participants is minimal: sessions consist of solving software-engineering tasks in a sandboxed browser workspace and answering Likert-scale post-task surveys, with no collection of sensitive identifiers beyond institutional email used for de-duplication. Participants are informed of the study purpose and data use at session start. The pilot is covered by an MIT COUHES exempt determination under 45 CFR 46.104(d)(2), Exempt Category 2. The protocol identifier is omitted from the submission and will be added in the camera-ready if required.
    \item[] Guidelines:
    \begin{itemize}
        \item The answer \answerNA{} means that the paper does not involve crowdsourcing nor research with human subjects.
        \item Depending on the country in which research is conducted, IRB approval (or equivalent) may be required for any human subjects research. If you obtained IRB approval, you should clearly state this in the paper. 
        \item We recognize that the procedures for this may vary significantly between institutions and locations, and we expect authors to adhere to the NeurIPS Code of Ethics and the guidelines for their institution. 
        \item For initial submissions, do not include any information that would break anonymity (if applicable), such as the institution conducting the review.
    \end{itemize}

\item {\bf Declaration of LLM usage}
    \item[] Question: Does the paper describe the usage of LLMs if they are an important component of the core methods?
    \item[] Answer: \answerYes{}
    \item[] Justification: LLMs are the system under test. Models evaluated are listed in Appendix~\ref{app:lb90-leaderboard} and span four families. Anthropic: Claude Haiku 4.5, Sonnet 4.6, Opus 4.7. Google: Gemini-3 Flash, Gemini-3.1 Flash Lite, Gemini-3.1 Pro, Gemma 4 31B, Gemini-2.5 Pro. OpenAI: GPT-5.4, GPT-5.4 Mini, gpt-oss-20b, gpt-oss-120b. Alibaba: Qwen 3 family. The Verifier role-mixing measurement explicitly evaluates each provider as the deciding LLM. LLMs were also used by the authors as writing aids, such as for grammar and clarity edits. All reported numbers, analyses, and conclusions were checked by the authors against the experimental logs.
\end{enumerate}
